\section{Stacks}
\textbf{FIFO} (First In First Out) data structure.

\begin{enumerate}
    \setlength\itemsep{0em}
    \item \textbf{push(x)}: Add element $x$ to the top of the stack.
    \item \textbf{pop()}: Remove and return the top element of the stack.
    \item \textbf{isEmpty()}: Return true if the stack is empty, false otherwise.
\end{enumerate}

\subsection{Implementations}
\subsubsection{Array-based implementation}
Make an array of size \verb|N| and a variable \verb|t| to keep track of the top index.

\begin{list}{}{}
    \setlength\itemsep{0em}
    \item \textbf{push(x)}: If \verb|t < N-1|, set \verb|t++| and \verb|arr[t] = x|.
    \item \textbf{pop()}: If \verb|t != -1|, return \verb|arr[t]| and \verb|t--|.
    \item \textbf{isEmpty()}: Return true if \verb|t < 0|, false otherwise.
\end{list}
\timeComplexity{1}
\spaceComplexity{N}

\subsubsection{Linked list-based implementation}
Dynamic data structure. More efficient than array-based implementation.

\begin{python}
class Node:
    def __init__(self, value):
        self.value = value
        self.next = None
class Stack:
    def __init__(self):
        self.top = None
    def push(self, x):
        new_node = Node(x)
        new_node.next = self.top
        self.top = new_node
    def pop(self):
        if self.top is None:
            return None
        value = self.top.value
        self.top = self.top.next
        return value
    def isEmpty(self):
        return self.top is None
\end{python}

\timeComplexity{1}
\spaceComplexity{n}

\subsubsection{Dynamic array-based implementation}
Make array \verb|A| size \verb|N|, array \verb|B| size \verb|2N| and variable \verb|t| for top index.
\emph{Push} adds to \verb|A| and copies 2 from \verb|A| to \verb|B|. When \verb|A| is full, set \verb|A = B| and make a new array \verb|B| of size \verb|4N|.
\emph{Pop} removes from \verb|A| and removes the same index from \verb|B|.
When \verb|A| $<$ \verb|N/4|, set \verb|B = A| and make a new array \verb|A| of size \verb|N/2|.

\timeComplexity{1}
\spaceComplexity{n} 

\section{Queues}
\textbf{LIFO} (Last In First Out) data structure.

\begin{enumerate}
    \setlength\itemsep{0em}
    \item \textbf{enqueue(x)}: Add element $x$ to the end of the queue.
    \item \textbf{dequeue()}: Remove and return the front element of the queue.
    \item \textbf{isEmpty()}: Return true if the queue is empty, false otherwise.
\end{enumerate}

\subsection{Implementations}
\subsubsection{Array-based implementation}
Make an array of size \verb|N| and two variables \verb|f| and \verb|b| to keep track of the front and back indices.
The queue snakes around the array using modulo.

\begin{list}{}{}
    \setlength\itemsep{0em}
    \item \textbf{enqueue(x)}: if \verb|(b+1) % N != f| set \verb|arr[b] = x| and \verb|b = (b+1) % N|.
    \item \textbf{dequeue()}: if \verb|f != b|, return \verb|arr[f]|, set \verb|f = (f+1) % N|.
    \item \textbf{isEmpty()}: Return true if \verb|f == b|, false otherwise.
\end{list}
\timeComplexity{1}
\spaceComplexity{N}

\subsubsection{Linked list-based implementation}
Dynamic data structure. More efficient than array-based implementation.
\begin{python}
class Node:
    def __init__(self, value):
        self.value = value
        self.next = None
class Queue:
    def __init__(self):
        self.front = None
        self.back = None
    def enqueue(self, x):
        new_node = Node(x)
        if self.back is None:
            self.front = self.back = new_node
        else:
            self.back.next = new_node
            self.back = new_node
    def dequeue(self):
        if self.front is None:
            return None
        value = self.front.value
        self.front = self.front.next
        if self.front is None:
            self.back = None
        return value
    def isEmpty(self):
        return self.front is None
\end{python}

\timeComplexity{1}
\spaceComplexity{n}

\subsubsection{Dynamic array-based implementation}
Can be done with multiple arrays like with stacks with same complexity.